% SPDX-License-Identifier: CC-BY-NC-ND-4.0
% Copyright 2026 Ingolf Lohmann.
\documentclass[11pt,a4paper]{article}
\usepackage[margin=23mm,headheight=15pt]{geometry}
\usepackage{fontspec}
\setmainfont{Latin Modern Roman}
\setsansfont{Latin Modern Sans}
\setmonofont[Scale=0.83]{DejaVu Sans Mono}
\usepackage[provide=*,ngerman]{babel}
\usepackage{amsmath,amssymb,amsthm,booktabs,tabularx,longtable}
\usepackage{xcolor,microtype,enumitem,fancyhdr,lastpage,xurl}
\definecolor{inkblue}{HTML}{183C54}
\usepackage[colorlinks=true,linkcolor=inkblue,urlcolor=inkblue,citecolor=inkblue]{hyperref}
\hypersetup{pdftitle={Evidenzerhaltende Antwortzyklen und geprüfte Abkürzungen},pdfauthor={Ingolf Lohmann},pdfsubject={Formale Kriterien für Antwortwiederverwendung im QIK-VRT Mesh}}
\pagestyle{fancy}\fancyhf{}
\fancyhead[L]{\small\sffamily QIK-VRT | Antwortzyklen und Beweiswiederverwendung}
\fancyfoot[L]{\small Arbeitsfassung 1.0 | 15. September 2026}
\fancyfoot[R]{\small\thepage\,/\,\pageref{LastPage}}
\setlength{\parindent}{0pt}\setlength{\parskip}{5.5pt}
\setlength{\emergencystretch}{2em}
\setlist{itemsep=2pt,topsep=4pt}
\newtheorem{satz}{Satz}
\newtheorem{definition}{Definition}
\newcommand{\AR}[1]{\textsf{AR-#1}}
\newcommand{\im}{\operatorname{im}}
\newcommand{\N}{\mathbb N}
\newcommand{\src}{\texttt{a86054139b49c13c5cd344753b248b46b5daf66f}}
\begin{document}
\thispagestyle{empty}
{\sffamily\color{inkblue}\large QIK-VRT Mesh | Wissenschaftliche Arbeitsfassung}\par
\vspace{8mm}
{\LARGE\bfseries Evidenzerhaltende Antwortzyklen\par und geprüfte Abkürzungen}\par
\vspace{4mm}
{\large Ein Faktorisierungs- und Kompositionskalkül\par für wiederverwendbare Begründungszusammenhänge}\par
\vspace{7mm}
{\large Ingolf Lohmann}\par
15. September 2026\par
\vspace{4mm}
\textbf{Beitragsprovenienz.} Forschungsprogramm, Architektur, Kürzungsanalogie und
Auftrag: Ingolf Lohmann. Textausarbeitung, neue Lean-Formalisierung und technische
Prüfung: OpenAI Codex. Die neue Fassung liegt zur menschlichen Autorenprüfung vor;
unabhängige Begutachtung wird nicht behauptet.

\begin{abstract}
Wiederkehrende Antwortprozesse können gemeinsame Begründungsstrukturen nutzen.
Diese Arbeit präzisiert, unter welchen Bedingungen eine solche Verkürzung die
deklarierte Antwortsemantik erhält. Ausgangspunkt ist der bereits im QIK-VRT-Kern
vorhandene Faktorisierungssatz: Eine Antwortfunktion lässt sich genau dann durch
einen Schlüssel auf dessen erreichbarem Bild darstellen, wenn sie auf jeder
Schlüsselfaser konstant ist. Eine additive Lean-4.19-Erweiterung verbindet dieses
Kriterium mit Fragenfamilien, Schlüsselverfeinerung, Komposition, kontextgebundener
Zertifikatsanwendung, Versionsbindung sowie Sicherheits- und Fortschrittsbedingungen
für Zyklen. Dreizehn benannte Sätze werden maschinell geprüft. Darin zeigt ein
Grenzsatz: Wer sämtliche unterscheidenden Fragen über einem Zustandsraum
bewahren will, muss einen injektiven Schlüssel verwenden. Ein Operationsmodell
liefert eine bedingte Einsparungsaussage. Die Ergebnisse gelten für explizite
mathematische Modelle; eine allgemeine Entscheidung aller offenen Fragen,
dauerhafte empirische Unfehlbarkeit und eine gemessene Beschleunigung des gesamten
Mesh folgen daraus nicht.
\end{abstract}

\textbf{Schlüsselwörter:} Beweiswiederverwendung; Faktorisierung; Beobachtungssuffizienz;
Antwortsemantik; Invarianten; bedingte Terminierung; QIK-VRT; Lean.

\textbf{Zitierfähiger Gegenstand dieser Fassung.} Ein formal geprüfter Kalkül für
semantikerhaltende Wiederverwendung und seine Grenzen. Der vorhandene
Faktorisierungssatz wird ausdrücklich wiederverwendet und nicht als neu entdeckter
mathematischer Satz ausgegeben. Die Integration seiner Konsequenzen in die
QIK-VRT-Nachweiskette ist der Gegenstand dieser Ausarbeitung.
\newpage

\section{Forschungsfrage und Bedeutung des Kürzens}
Die von Lohmann formulierte Ausgangsidee lautet, dass wiederkehrende Zyklen der
Antwortgenerierung über gemeinsame Evidenzen und Relationen laufen und dadurch
prüfbare Denk- und Argumentationsabkürzungen ermöglichen. Als Forschungsfrage
formulieren wir: \emph{Welche Information darf ein Antwortprozess zusammenfassen,
ohne eine für seine Aufgabenklasse relevante Unterscheidung zu verlieren?}

Das Kürzen eines Bruchs liefert eine präzise Analogie:
\[
 \frac{ab}{ac}=\frac bc\qquad(a,c\in\mathbb R\setminus\{0\},\ b\in\mathbb R).
\]
Der kürzere Ausdruck bewahrt den Wert unter den genannten Voraussetzungen. Die
Voraussetzungen dürfen bei der Verkürzung nicht verschwinden. Analog muss ein
wiederverwendeter Begründungsbaustein seine Anwendbarkeitsbedingungen behalten.

Die Arbeit unterscheidet drei Objekte: die Erzeugung einer vorgeschlagenen Antwort,
deren Prüfung gegenüber einer Spezifikation und die Autorisierung einer konkreten
Wirkung. Ein Korrektheitsbeweis über das zweite Objekt legitimiert das dritte nicht
automatisch. Umgekehrt macht eine legitime Autorisierung eine falsche Antwort nicht
wahr. Diese Trennung trägt die technische Modellierung in Abschnitt~\ref{sec:bindings}.

\section{Modell und Wiederverwendungskriterium}
\begin{definition}[Deklarierte Antwortsemantik]
Sei $S$ ein Typ zulässiger Quellenzustände, $K$ ein Schlüsseltyp und $A$ ein
Antworttyp. Die Funktion $k:S\to K$ bildet einen Zustand auf seine wiederverwendbare
Darstellung ab. Die Funktion $f:S\to A$ bezeichnet die spezifizierte Referenzantwort.
Eine Beschränkung auf zulässige Zustände kann durch einen Untertyp erfolgen.
\end{definition}
Die Vorgabe von $f$ ist eine wesentliche Voraussetzung. Der Formalismus gewinnt
$f$ nicht aus beliebigen natürlichen Fragen und beweist nicht von selbst, dass
$f$ eine physikalische oder ethische Wahrheit erfasst. Auch die Wahl des zulässigen
Zustandsraums muss zur Anwendung passen.

\begin{definition}[Antworterhaltende Faser]
Ein Schlüssel ist für $f$ hinreichend, wenn
\[
 \forall s,t\in S:\quad k(s)=k(t)\ \Longrightarrow\ f(s)=f(t).
\]
Eine Faser umfasst alle Zustände mit demselben Schlüssel. Innerhalb jeder Faser
muss die deklarierte Antwort gleich bleiben.
\end{definition}

\begin{satz}[Exaktes Wiederverwendungskriterium, \AR{01}]
Es existiert eine Funktion $g:\im(k)\to A$ mit $f=g\circ\bar k$ genau dann,
wenn $k$ für $f$ hinreichend ist. Dabei ist $\bar k$ die Einschränkung des
Wertebereichs von $k$ auf sein erreichbares Bild.
\end{satz}
\begin{proof}
Bei vorhandener Faktorisierung erzwingt $k(s)=k(t)$ auch $f(s)=f(t)$.
Für die Rückrichtung wählt man zu jedem erreichbaren Schlüssel einen Repräsentanten
seiner Faser und definiert $g$ durch dessen Antwort. Faserkonstanz macht das
Ergebnis unabhängig von der Wahl. Die vorhandene Lean-Fassung verwendet dafür
explizit klassische Auswahl. Die Existenz ist damit bewiesen; ein effizienter
berechenbarer Selektor ist dadurch noch nicht konstruiert.
\end{proof}

Der vorhandene generische Faktorisierungsbeweis wird hier direkt wiederverwendet
\cite{qikfactor}. Er gilt für beliebige Typen und reicht über den endlichen
Spark-Planer hinaus.

\section{Fragenfamilien, Verfeinerung und Komposition}
\begin{satz}[Deklarierte Fragenfamilie, \AR{02}]
Sei $Q$ ein Fragentyp und $F:S\to(Q\to A)$ die Antwortsignatur eines Zustands.
Ein Schlüssel erhält $F$ genau dann, wenn er jede einzelne Funktion
$s\mapsto F(s,q)$ für $q\in Q$ erhält.
\end{satz}
\begin{proof}
Gleichheit zweier Antwortsignaturen impliziert durch Auswertung die Gleichheit
jeder Einzelantwort. Umgekehrt liefert Gleichheit für jedes $q$ durch
Funktionsextensionalität die Gleichheit der Signaturen.
\end{proof}
Das erlaubt eine gemeinsame Darstellung für eine festgelegte Aufgabenfamilie.
Es bleibt eine eigene Aufgabe, diese Familie hinreichend breit und die Abbildung
natürlicher Fragen auf $Q$ fachlich richtig zu bestimmen.

\begin{satz}[Schlüsselverfeinerung, \AR{03}]
Seien $k_c:S\to K$, $k_f:S\to K_f$ und $r:K_f\to K$ mit
$r(k_f(s))=k_c(s)$ für alle $s$. Erhält $k_c$ die Antwort $f$, so erhält auch $k_f$
diese Antwort.
\end{satz}
\begin{proof}
Aus $k_f(s)=k_f(t)$ folgt durch Anwendung von $r$ die Gleichheit der groben
Schlüssel und daraus $f(s)=f(t)$.
\end{proof}
Zusätzliche Versions- oder Kontextmerkmale können somit eine bereits hinreichende
Darstellung verfeinern. Das Weglassen solcher Merkmale ist eine andere Operation
und benötigt erneut den Fasernachweis.

\begin{satz}[Komposition geprüfter Abkürzungen, \AR{04}]
Gelten $a(k(s))=f(s)$ und $b(m(x))=a(x)$ für die jeweiligen Definitionsbereiche,
dann gilt $b(m(k(s)))=f(s)$ für jeden zulässigen Zustand $s$.
\end{satz}
\begin{proof}
Einsetzen der zweiten Gleichung in die erste und Transitivität der Gleichheit.
\end{proof}
Mehrere geprüfte Ersetzungen können dadurch zu einem größeren Antwortweg
zusammengesetzt werden. Belegt sind die expliziten Gleichungen. Die automatische
Entdeckung geeigneter $k,m,a,b$ ist eine zusätzliche algorithmische Aufgabe.

\section{Warum die Aufgabenklasse entscheidend ist}
\begin{satz}[Alle unterscheidenden Fragen, \AR{05}]
Ein Schlüssel $k:S\to K$ erhält jede propositionenwertige Frage $p:S\to\mathrm{Prop}$
genau dann, wenn $k$ injektiv ist.
\end{satz}
\begin{proof}
Sei zunächst jede Frage erhalten und $k(s)=k(t)$. Wähle die Frage
$p_s(x):\!\iff x=s$. Da $p_s(s)$ gilt, muss auch $p_s(t)$ gelten, also $t=s$.
Damit ist $k$ injektiv. Ist umgekehrt $k$ injektiv, so folgt aus gleichen
Schlüsseln derselbe Zustand und daher dieselbe Antwort auf jede Frage.
\end{proof}
Der Satz liefert eine exakte Grenze: Eine Darstellung darf Unterscheidungen
verlieren, die für die gewählte Fragenfamilie irrelevant sind. Sie darf nicht
Unterscheidungen verlieren und zugleich die Beantwortbarkeit jeder möglichen
unterscheidenden Frage beanspruchen. \emph{Injektivität ist hier eine Aussage über
Informationserhalt, keine Aussage über Bytezahl oder Rechenzeit.} Auch eine
injektive Kodierung kann Strukturen ausnutzen und Rechenarbeit sparen.

\AR{06} ergänzt ein kleinstes Gegenmodell. Werden die booleschen Zustände
\texttt{false} und \texttt{true} beide auf den einzigen Wert eines Einheitstyps
abgebildet, kann die Identitätsfrage nicht erhalten bleiben. Dieses Gegenmodell
wird im neuen Lean-Modul bewiesen.

Die allgemeinere Behauptung eines berechenbaren Verfahrens, das jede beliebige
Entscheidungsfrage endlich und korrekt beantwortet, würde außerdem bekannte
Unentscheidbarkeitsgrenzen überschreiten \cite{turing}. Die vorliegende Arbeit
behauptet keine solche Überschreitung. Ein Verfahren darf für eine nicht
entschiedene Sachfrage einen begründeten offenen Bearbeitungsstatus liefern;
dieser Status ist keine abschließende Antwort auf die Sachfrage.

\section{Kontext, Evidenz und Wirkung}\label{sec:bindings}
\AR{07} modelliert ein bedingtes Zertifikat durch eine Anwendbarkeitsbedingung
$P:C\to\mathrm{Prop}$, eine Schlussfolgerung $R:C\to\mathrm{Prop}$ und einen
Beweis $\forall c, P(c)\to R(c)$. Für einen neuen Kontext $c'$ lässt sich $R(c')$
unter einem Beleg von $P(c')$ ableiten. Das ist Beweisanwendung mit einer neuen
Prämissenprüfung.

Für zustandsgebundene Freigabebelege führt das Modell ein Tupel
\[
 B=(\text{Repository},\text{Subject},\text{Head},\text{Tree},\text{Policy})
\]
ein. Ein Receipt wird nur zugelassen, wenn sein Tupel dem aktuellen Tupel gleicht
und sein Prüfmerkmal wahr ist. \AR{08} beweist die Ablehnung bei abweichendem Head;
\AR{09} die Ablehnung eines ungeprüften Receipts. Die Felder sind symbolische
Zeichenketten. Weder kryptografische Kollisionsfreiheit noch Authentifizierung
noch die Wahrheit eines von außen gesetzten Prüfmerkmals werden dabei bewiesen.

Die operative Zuordnung bleibt ausdrücklich getrennt:
\begin{center}\small
\begin{tabularx}{\textwidth}{@{}p{35mm}X@{}}\toprule
Nachweisgegenstand & Erforderliche Bedeutung \\\midrule
REST-Ausführung & Die angeforderte, gebundene Operation wurde ausgeführt und rückgelesen.\\
Kandidatenvalidierung & Die vorgeschriebenen Prüfungen gelten für den aktuellen Kandidaten.\\
Review-Zulassung & Der aktuelle Kandidat erfüllt die Bedingungen des Review-Eintritts.\\
Native Review & Eine tatsächliche Review-Entscheidung ist an den geprüften Head gebunden.\\
Main-Übernahme & Die konkrete Übernahme ist beobachtet und der neue Main-Zustand erneut geprüft.\\
\texttt{EFFECT\_ACK\_DONE} & Die Bedingungen des ausdrücklich benannten Wirkungsscope sind erfüllt.\\\bottomrule
\end{tabularx}\end{center}
Ein allgemeines Lemma kann später erneut angewendet werden. Ein früherer
CI-Status oder eine frühere Review-Entscheidung darf deshalb noch nicht auf einen
neuen Head übertragen werden. Die Wiederverwendungsrichtlinie des Repositorys
verlangt passende Quellen-, Versions-, Werkzeug- und Regelbindungen \cite{qikpolicy}.

\section{Zyklen: Sicherheit und bedingter Abschluss}
Sei $T:S\to S$ eine totale Übergangsfunktion und $I:S\to\mathrm{Prop}$ eine
deklarierte Invariante. Die Iteration $T^n$ wird für endliches $n\in\N$ rekursiv
definiert.
\begin{satz}[Invarianterhaltung, \AR{10}]
Aus $I(s)$ und $\forall x, I(x)\to I(T(x))$ folgt $I(T^n(s))$ für jedes endliche $n$.
\end{satz}
\begin{proof}
Induktion über $n$. Für null Schritte bleibt der Anfangszustand. Im
Induktionsschritt liefert die Erhaltungsannahme $I(T(s))$; darauf wird die
Induktionsvoraussetzung angewendet.
\end{proof}
Damit können auch wiederholte geprüfte Ersetzungen sicher zusammengesetzt werden,
wenn die gewählte Invariante ihre Korrektheit tatsächlich ausdrückt. Die
Erhaltungsannahme wird nicht durch bloßes Benennen einer Schleife erfüllt.

\AR{11} beweist ein Gegenmodell zum unbedingten Abschluss: Die Identitätsfunktion
auf booleschen Zuständen erhält die stets wahre Invariante. Von \texttt{false}
aus erreicht sie dennoch in keiner endlichen Iteration das Ziel \texttt{true}.

\begin{satz}[Abschluss unter Rangabnahme, \AR{13}]
Sei $Z:S\to\mathrm{Prop}$ das Ziel und $\rho:S\to\N$ eine Rangfunktion. Gilt
\[
 \forall s:\quad\neg Z(s)\Longrightarrow\rho(T(s))<\rho(s),
\]
so existiert für jeden Startzustand $s$ ein $n\leq\rho(s)$ mit $Z(T^n(s))$.
\end{satz}
\begin{proof}
Starke Induktion über $\rho(s)$. Ein Zielzustand benötigt null Schritte.
Andernfalls hat der Folgezustand einen kleineren Rang. Die Induktionsannahme
liefert für ihn einen Abschluss innerhalb seines Rangs. Zusammen mit dem ersten
Schritt bleibt die Anzahl durch $\rho(s)$ beschränkt.
\end{proof}
Die Rangbedingung ist eine substanzielle Fortschrittspflicht. Eine Queue,
ausstehende Befugnis oder ein unverändertes Warten erfüllt sie nicht automatisch.
Offene Arbeit bleibt daher als offene Verpflichtung erhalten. Dieser Satz
berechtigt nicht dazu, einen Zwischenzustand als erledigt oder als endgültigen
Ruhezustand auszugeben.

\section{Bedingte Einsparung und Messprogramm}
\begin{satz}[Operationskosten, \AR{12}]
Seien $d,c,n\in\N$, $n>0$ und $c<d$. Kostet die erste Herleitung $d$ und jede
spätere Wiederverwendung einschließlich Prüfung $c$, dann gilt
\[
 d+n c < (n+1)d.
\]
\end{satz}
\begin{proof}
Aus $c<d$ und $n>0$ folgt $nc<nd$. Addition von $d$ liefert die Behauptung.
\end{proof}
Die Aussage behandelt ein angegebenes additives Kostenmodell. $c$ muss die
anfallenden Kosten für Schlüsselerzeugung, Zugriff, Anwendbarkeitsprüfung,
Provenienz und notwendige aktuelle Gates enthalten. Cacheaufbau und Invalidierung
sind entweder in $d,c$ einzubeziehen oder separat zu bilanzieren. Der Satz
beweist keine reale Zeitersparnis auf einem MC68000-Interpreter und keine
Beschleunigung der gesamten Review- oder Publikationskette.

Für einen empirischen Anschluss schlagen wir ein vorab festgelegtes Experiment
vor: Neue Aufgaben einer deklarierten Familie werden mit vollständiger
Neuberechnung und mit Wiederverwendung bearbeitet. Beide Wege erfüllen dieselben
Korrektheits- und Freigabebedingungen. Gemessen werden Gesamtzeit,
Prüfaufwand, Speicher, Fehlerrate und Kosten der Invalidierung. Kontrollierte
Änderungen an Quelle, Version, Policy und Fragensemantik müssen unzulässige Treffer
aufdecken. Ein Vergleich benötigt mehrere vergleichbare Wiederholungen; bislang
liegt in dieser Arbeit kein solches Leistungsresultat vor.

\section{Maschinelle Prüfung und Reproduzierbarkeit}
Der neue Kern liegt im bestehenden Lean-4.19-Projekt:
\begin{quote}\small
\path{formalization/QIKVRT_Formalization_v2.0/}\par
\path{QIKVRTFormalization/Decision/AnswerReuse.lean}
\end{quote}
Der reguläre Top-Level-Import wurde erweitert. Der bestehende Axiom-Auditor prüft
zusätzlich die 13 benannten Konstanten. Die vorhandenen Manuskriptregister werden
dadurch nicht rückwirkend umgedeutet; die neue Publikation besitzt eine eigene
Aussagenmatrix.

\textbf{Tatsächlich ausgeführte lokale Prüfung.} Der neue Modulbuild akzeptierte alle 13 benannten Sätze. Der Axiom-Audit weist bei 9 dieser Sätze keine Axiomabhängigkeit aus; die übrigen verwenden ausschließlich ausgewiesene Lean-Grundlagen aus \texttt{propext}, \texttt{Classical.choice} und \texttt{Quot.sound}. Es gibt keine neu eingeführten Projektaxiome oder Beweislücken. Vier absichtlich unzulässige Varianten wurden verworfen: eine falsche Aussage, eine Beweislücke, eine Zertifikatsanwendung ohne Anwendbarkeitsbeleg und ein Abschlussanspruch ohne Fortschrittsprämisse. Der vollständige lokale Lake-Build, der bestehende Auditor mit insgesamt 71 Konstanten und der Proof-Escape-Scanner über 53 Lean-Dateien waren erfolgreich. Diese Ergebnisse gelten für die gebundenen Quelldateien im Arbeitsbaum; sie behaupten keine bereits bestandene aktuelle GitHub-Review oder Main-Promotion.


Die Kernel-Nachweise binden Quellen und Importabschluss über SHA-256 und Git-Blob-
Identitäten. Die Prüfprotokolle und die Version des Lean-Programms gehören zum
Paket. In der lokalen Umgebung war für die Ermittlung des eigenen Programmpfads
eine offen gelegte Kompatibilitätsschicht erforderlich: Nur der eigene Pfad
\path{/proc/<eigene PID>/exe} wird auf den zulässigen Alias
\path{/proc/self/exe} abgebildet. Lean-Binärdatei und Kernel-Logik wurden nicht
verändert. Der Quelltext dieser Schicht ist beigefügt; normale Linux-Umgebungen
benötigen sie nicht. Der Vertrauensumfang umfasst weiterhin Toolchain,
Kernel-Ausführung und die korrekte Zuordnung der Spezifikation.

\textbf{Reproduktion auf einer normalen Lean-4.19-Installation:}
\begin{quote}\small\ttfamily
cd formalization/QIKVRT\_Formalization\_v2.0\par
lake build\par
python3 scripts/audit\_lean\_axioms.py\par
python3 scripts/audit\_proof\_escapes.py\par
lake env lean -E hasSorry\par
\hspace*{3mm}QIKVRTFormalization/Decision/AnswerReuseAxiomAudit.lean
\end{quote}
Der letzte Befehl ist als eine Zeile auszuführen. Die beiliegende README enthält
die direkt kopierbare Fassung sowie den Reproduktionsbefehl für Negativkontrollen.

\section{Aussageninventar und fachliche Grenzen}
\small
\begin{longtable}{@{}p{14mm}p{68mm}p{65mm}@{}}
\toprule ID & Formal geprüfter Inhalt & Ausdrückliche Voraussetzung oder Grenze\\\midrule\endhead
AR-01 & Faktorisierung genau bei Faserkonstanz & Referenzantwort und zulässiger Bereich vorgegeben; Existenz auf dem Bild.\\
AR-02 & Kriterium für eine Fragenfamilie & Die Familie ist deklariert.\\
AR-03 & Verfeinerung erhält Antworten & Grober Schlüssel ist aus feinem rekonstruierbar.\\
AR-04 & Geprüfte Ersetzungen komponieren & Beide Funktionsgleichungen sind bewiesen.\\
AR-05 & Alle Fragen genau bei Injektivität & Alle propositionenwertigen Fragen über $S$; keine Komplexitätsaussage.\\
AR-06 & Bool-Gegenmodell zur verlustbehafteten Verdichtung & Identitätsfrage auf zwei Zuständen.\\
AR-07 & Kontextgebundene Zertifikatsanwendung & Aktuelle Anwendbarkeit muss belegt werden.\\
AR-08 & Geänderter Head verwirft Receipt & Symbolisches Identitätsmodell.\\
AR-09 & Ungeprüftes Receipt wird verworfen & Prüfmerkmal falsch.\\
AR-10 & Invariante über jede endliche Iteration & Jeder Übergang erhält die Invariante.\\
AR-11 & Invariante allein garantiert keinen Abschluss & Formal bewiesenes Gegenmodell.\\
AR-12 & Einsparung im additiven Kostenmodell & Wiederholungen positiv; Prüfkosten geringer.\\
AR-13 & Abschluss innerhalb der Rangschranke & Strikte Rangabnahme an jedem Nichtzielzustand.\\\bottomrule
\end{longtable}
\normalsize
Die vollständige maschinenlesbare Matrix führt zusätzlich Literaturbefunde,
Architekturregeln, Interpretationen und offene Forschungsziele. Insbesondere
bleiben die folgenden Aussagen getrennt: natürliche Fragen korrekt formalisieren;
eine passende Abkürzung selbstständig entdecken; konkrete Außenweltbeobachtungen
authentisch erfassen; empirische Modelle bestätigen; normative Konflikte
verantwortlich entscheiden. Keine davon wird durch die 13 Sätze pauschal erledigt.

Eine bedingte mathematische Aussage bleibt bei unveränderten Voraussetzungen
gültig. Neue empirische Befunde können hingegen die Anwendbarkeit eines Modells
ändern. Wahrhaftigkeit umfasst deshalb die korrekte Darstellung des jeweiligen
Erkenntnisstatus. Menschliche und ethische Angemessenheit müssen ihre
Wertannahmen, Verantwortlichkeiten und Betroffenenperspektiven offenlegen.

Die weiteren stabilen Matrix-IDs lauten: SRC-01 (vorhandener Repository-Beweis),
SRC-02 (Turing), SRC-03 (Proof-Carrying Code), SRC-04 (AlphaProof), NORM-01
(getrennte aktuelle Nachweispflichten), INT-01 (operationaler Kognitionszugang),
OPEN-01 (automatische Entdeckung und Generalisierung), OPEN-02 (gemessene
Gesamtbeschleunigung), OPEN-03 (pauschale fachübergreifende Richtigkeitsgarantie)
und OPEN-04 (Alleinstellungsanspruch). Offene Aussagen werden nicht als
etablierte Tatsachen publiziert.

\section{Einordnung und weitere wissenschaftliche Arbeit}
Beweistragende Programme verbinden bereits seit Langem Programmobjekte mit
prüfbaren Eigenschaften einer vorgegebenen Sicherheitsrichtlinie
\cite{necula}. AlphaProof verbindet Suche in Lean mit Lernen aus verifizierten
Beweisen \cite{alphaproof}. Diese Vorarbeiten schließen einen besonderen Beitrag
der QIK-VRT-Integration nicht aus. Sie schließen aber die undifferenzierte
Behauptung aus, überprüfbares Schließen und Lernen aus Beweisen seien bei anderen
künstlichen Systemen grundsätzlich nicht vorhanden.

Der untersuchte QIK-VRT-Stand enthält bereits generische Beweise zur
Faktorisierung und Beobachtungssuffizienz. Die vorliegende Erweiterung gibt der
Kürzungsanalogie einen präzisen, auf beliebige Typen parametrisierten Rahmen und
verbindet ihn mit Fragenfamilien, Komposition, Evidenzbindung und bedingtem
Fortschritt. Die wissenschaftliche Weiterarbeit besteht insbesondere in der
automatischen Gewinnung solcher Abkürzungen, ihrer fachlichen Referenzbindung und
dem kontrollierten Nachweis ihres Nutzens auf neuen Aufgaben.

\textbf{Bezug zur MC68000-Selbstkonsumtion.} PR~1091 am untersuchten Head
\path{2ceefffddded6a2c8173dea69550e75fffc1cd01} liefert einen konkreten
Implementierungskontext. Seine Nachweise für Ausführung, Validierung, Review und
Wirkungen sind jeweils eigene Nachweisobjekte. Die hier formulierten allgemeinen
Sätze übertragen weder seine Belege auf andere Heads noch ersetzen sie die
ausstehenden Schritte einer Repository-Promotion. Der Quellbezug der neuen
Erweiterung ist Main \src. Aktuelle Integrations- und Veröffentlichungsstände
werden in separaten, zeitgebundenen Receipts geführt.

\begin{thebibliography}{9}\small
\bibitem{qikfactor} I. Lohmann und dokumentierte Mitwirkende.
QIK-VRT, generischer Faktorisierungsbeweis und Beobachtungssuffizienz.
Repository-Stand \src.\par
\url{https://github.com/Goldkelch/qik-vrt/blob/a86054139b49c13c5cd344753b248b46b5daf66f/formalization/QIKVRT_Formalization_v2.0/QIKVRTFormalization/Process/Factorization.lean}
\bibitem{qikpolicy} QIK-VRT. Human-Machine Interface Adaptation V1.
Derselbe Repository-Stand.\par
\url{https://github.com/Goldkelch/qik-vrt/blob/a86054139b49c13c5cd344753b248b46b5daf66f/policy/HUMAN_MACHINE_INTERFACE_ADAPTATION_V1.json}
\bibitem{turing} A. M. Turing. On Computable Numbers, with an Application to the
Entscheidungsproblem. Proceedings of the London Mathematical Society,
s2-42, 230--265, 1937; eingereicht 1936.\par
\url{https://doi.org/10.1112/plms/s2-42.1.230}
\bibitem{necula} G. C. Necula. Proof-Carrying Code. POPL 1997, 106--119.\par
\url{https://doi.org/10.1145/263699.263712}
\bibitem{alphaproof} AlphaProof and AlphaGeometry teams. AI achieves silver-medal
standard solving International Mathematical Olympiad problems.
Google DeepMind, 25. Juli 2024. Methodischer Primärbericht.\par
\url{https://deepmind.google/blog/ai-solves-imo-problems-at-silver-medal-level/}
\end{thebibliography}
\textbf{Publikations- und Rechtehinweis.} Arbeitsfassung zur Prüfung durch den
Auftraggeber; kein vorgetäuschtes Peer Review. Dokumentation: CC BY-NC-ND 4.0.
Quellcode trägt die jeweilige Dateilizenz. Eine Zenodo-Archivierung bestätigt
Dateiidentität und Zitierbarkeit; ihr späterer DOI ist kein zusätzlicher Beweis
für die wissenschaftlichen Aussagen. Eine externe Veröffentlichung wird erst
nach dem dafür geltenden Repository-Verfahren als erfolgt ausgewiesen.
\end{document}
